\subsection{A genuine higher-height scalar kernel test}\label{sec:kerneltest}

The unchanged-logarithm obstruction DP21 leaves a concrete alternative: replace the logarithm itself and test the actual analytic hypotheses. Fix an odd prime, $h\ge2$, $q=p^h$, and put
\[
 L_h(T)=\sum_{e\ge0}T^{q^e}/p^e,\qquad E_h=L_h^{-1},\qquad
 \rho_h=(q-1)^{-1},\quad s=(p-1)^{-1}.
\]
The Honda formal group $F_h(S,T)=E_h(L_h(S)+L_h(T))$ is integral and has height $h$: integrality is the functional-equation lemma of Hazewinkel \cite[Lemma 2.3 and (2.5), p.~74]{Hazewinkel}, with coefficient ring $\Z_p$, $\sigma=\id$, and first functional-equation coefficient $1/p$. Its $p$-series begins $pT+(1-p^{q-1})T^q$. DP2, or direct scaling of this logarithm, gives
\begin{equation}\label{eq:honda-scale}
 E_{h,\pi}(Z)=\pi^{-1}E_h(\pi Z)\in R'[[Z]],\qquad
 \overline E_{h,\pi}(Z)=\sum_{r\ge0}(-1)^rZ^{q^r},\quad \pi^{q-1}=p.
\end{equation}
Here and below $R'$ denotes the valuation ring after adjoining $\pi$. In particular $v([T^n]E_h)\ge-\rho_h(n-1)$, with equality for $n=q^r$.

For this test, let $K$ be a valued extension of $\Q_p$, with $v(p)=1$, containing the $(q-1)$-st roots of unity, and let $R$ be its valuation ring, and choose a primitive root $\zeta$. For $\beta\in K[[X,Y]]$ impose the following explicitly modified scalar hypotheses:
\begin{enumerate}[label=(H\arabic*),leftmargin=2.8em]
\item its coefficients lie in an unramified subfield;
\item $\beta(\zeta X,\zeta^{-p}Y)=\beta(X,Y)$;
\item every nonzero monomial has positive bidegrees, and $\beta_{1,1}=0$;
\item the two addition coboundaries
\[
 H_X=\exp\bigl(\delta_X\beta(L_h(S),L_h(S'),L_h(T))\bigr),\qquad
 H_Y=\exp\bigl(\delta_Y\beta(L_h(S),L_h(T),L_h(T'))\bigr)
\]
are integral units with constant coefficient one. Here $\delta_X\beta=\beta(X+X',Y)-\beta(X,Y)-\beta(X',Y)$, and similarly for $Y$.
\end{enumerate}
This keeps the characters $1,-p$, replaces the root order by $p^h-1$, and replaces both scalar logarithms. It is a well-defined analytic proposal, not the assertion that a one-dimensional height-$h$ group is its own Cartier dual or that it realizes an $h$-letter theta correspondence. At $h=2$ these are the scalar conditions in \path{GaussJetBounds/Statement.lean}, \path{BHyp}, lines 39--94.

